\section{Kategorie triangulowane i pochodne}\label{cat:sec:triangulated-derived}
W tym rozdziale przechodzimy od pojedynczych grup homologii i funktorów pochodnych do kategorii, w których kompleksy są obiektami, a homotopie i quasi-izomorfizmy wpływają na pojęcie morfizmu. Korzystamy z rozdziałów~\ref{cat:sec:abelian-categories}, \ref{cat:sec:projective-injective} i~\ref{cat:sec:derived-functors}. Wersję łańcuchową homotopii omawia \atlasref{TOP-03::def:homotopia-lancuchowa}{TOP-03: homotopia łańcuchowa}. Przyjmujemy w całym rozdziale konwencję kołańcuchową.

\subsection{Kategoria homotopijna kompleksów}
\begin{definicja}\label{cat:def:complex-category}
Dla kategorii addytywnej $\mathcal A$ przez $\cat{C}(\mathcal A)$ oznaczamy kategorię kompleksów kołańcuchowych $X=(X^n,d_X^n)_{n\in\mathbb Z}$ z $d_X^{n+1}d_X^n=0$. Morfizm $f:X\to Y$ jest rodziną $f^n:X^n\to Y^n$ spełniającą $d_Y^nf^n=f^{n+1}d_X^n$. Homotopią między $f,g:X\to Y$ nazywamy rodzinę $h^n:X^n\to Y^{n-1}$ taką, że
\begin{equation}\label{cat:eq:cochain-homotopy}
 f^n-g^n=d_Y^{n-1}h^n+h^{n+1}d_X^n.
\end{equation}
Piszemy wówczas $f\simeq g$. Kategoria $\cat{K}(\mathcal A)$ ma te same obiekty co $\cat{C}(\mathcal A)$, a morfizmy są klasami homotopii. Symbole $\cat{C}^+,\cat{C}^-,\cat{C}^b$ i analogicznie $\cat{K}^+,\cat{K}^-,\cat{K}^b$ oznaczają pełne podkategorie kompleksów ograniczonych odpowiednio z dołu, z góry i z obu stron (w sensie zerowania się obiektów poza odpowiednim zakresem stopni).
\end{definicja}
\begin{stwierdzenie}\label{cat:prop:homotopy-category}
Kategoria $\cat{K}(\mathcal A)$ jest dobrze określoną kategorią addytywną, a funktor ilorazowy $\cat{C}(\mathcal A)\to\cat{K}(\mathcal A)$ jest addytywny.
\end{stwierdzenie}
\begin{proof}
Morfizmy kompleksów tworzą grupę abelową z dodawaniem stopniowym. Dla dowolnej rodziny $h^n:X^n\to Y^{n-1}$ morfizmy $u^n=d_Y^{n-1}h^n+h^{n+1}d_X^n$ komutują z różniczkami: obie strony równania $d_Y^nu^n=u^{n+1}d_X^n$ są równe $d_Y^nh^{n+1}d_X^n$. Takie morfizmy tworzą podgrupę, ponieważ suma i przeciwieństwo rodzin $h$ realizują odpowiednio sumę i przeciwieństwo $u$. Relacja homotopii jest więc relacją równoważności i iloraz zbioru morfizmów jest grupą abelową.
Jeśli $f-g=d_Yh+hd_X$, $a:W\to X$ i $b:Y\to Z$ są morfizmami kompleksów, to
\[
 b(f-g)a=d_Z(bha)+(bha)d_W.
\]
Wzór wynika z przemienności $a,b$ z różniczkami, z zachowaniem odpowiednich stopni. Homotopijność jest zatem zgodna ze składaniem i składanie klas jest dobrze określone, dwuliniowe i łączne. Identyczności reprezentują identyczności kompleksów. Obiekt zerowy i biprodukty w kategorii kompleksów otrzymujemy stopniowo; ich równania definiujące zachodzą także po przejściu do klas, więc są biproduktami w kategorii homotopijnej. Otrzymujemy kategorię addytywną oraz addytywny funktor ilorazowy.
\end{proof}
\begin{definicja}\label{cat:def:contractible-complex}
Kompleks $X$ nazywamy ściągalnym, jeśli $\id_X\simeq0$. Morfizm $f:X\to Y$ jest równoważnością homotopijną, jeśli istnieje $g:Y\to X$ z $gf\simeq\id_X$ i $fg\simeq\id_Y$.
\end{definicja}
\begin{stwierdzenie}\label{cat:prop:contractible-zero}
Kompleks jest ściągalny wtedy i tylko wtedy, gdy jest obiektem zerowym w $\cat{K}(\mathcal A)$. Równoważności homotopijne to dokładnie morfizmy, których klasy są izomorfizmami. Jeśli $\mathcal A$ jest abelowa, każdy ściągalny kompleks jest acykliczny, tj. $H^n(X)=0$ dla wszystkich $n$.
\end{stwierdzenie}
\begin{proof}
W kategorii addytywnej obiekt jest zerowy dokładnie wtedy, gdy jego identyczność jest zerowa: każdy morfizm do niego i z niego jest złożeniem z tą identycznością, a więc zerem. W $\cat{K}(\mathcal A)$ oznacza to $\id_X\simeq0$. Definicja izomorfizmu klas daje wprost oba równania homotopijne dla reprezentanta odwrotności. W kategorii abelowej kohomologia zachowuje identyczności i utożsamia morfizmy homotopijne na mocy twierdzenia~\ref{cat:thm:abelian-homology-functor}. Stąd $\id_{H^n(X)}=H^n(\id_X)=H^n(0)=0$, więc $H^n(X)$ jest zerowe.
\end{proof}
\begin{przyklad}\label{cat:ex:acyclic-not-contractible}
Kompleks grup abelowych skoncentrowany w stopniach $-1,0,1$,
\[
 \mathbb Z\xrightarrow{m}\mathbb Z\longrightarrow\mathbb Z/m\mathbb Z
 \qquad(m\ge2),
\]
gdzie druga strzałka jest projekcją, jest acykliczny, ponieważ wraz z zerami na końcach jest ciągiem dokładnym. Nie jest jednak ściągalny. Gdyby $\id_X=d_Xh+hd_X$, w stopniu $1$ mielibyśmy $\id_{\mathbb Z/m\mathbb Z}=d_X^0h^1$, ponieważ $d_X^1=0$. Każdy homomorfizm $h^1:\mathbb Z/m\mathbb Z\to\mathbb Z$ jest zerowy (obraz generatora jest elementem torsyjnym grupy bez torsji), co daje sprzeczność. Por. \atlasref{TOP-03::prz:acykliczny-niesciagalny}{TOP-03: kompleks acykliczny, który nie jest ściągalny}.
\end{przyklad}

\subsection{Przesunięcie kompleksu}
\begin{definicja}\label{cat:def:complex-shift}
Dla $r\in\mathbb Z$ definiujemy kompleks $X[r]$ i morfizm $f[r]$ przez
\[
 X[r]^n=X^{n+r},\qquad d_{X[r]}^n=(-1)^r d_X^{n+r},
 \qquad f[r]^n=f^{n+r}.
\]
W szczególności $X[1]$ ma różniczkę przeciwną do odpowiednio przesuniętej różniczki $X$.
\end{definicja}
\begin{stwierdzenie}\label{cat:prop:shift-functor}
Przesunięcie $[r]$ wyznacza addytywny automorfizm $\cat{C}(\mathcal A)$ oraz $\cat{K}(\mathcal A)$, z odwrotnością $[-r]$. Zachodzą równości $[r][s]=[r+s]$ i $[0]=\id$. Jeżeli $\mathcal A$ jest abelowa, to $H^n(X[r])\cong H^{n+r}(X)$ naturalnie w $X$.
\end{stwierdzenie}
\begin{proof}
Kwadrat różniczki przesuniętej jest zerem, ponieważ mnożymy złożenie dwóch kolejnych różniczek przez $(-1)^{2r}=1$. Równanie morfizmu kompleksów zachowuje się po pomnożeniu przez $(-1)^r$ i zmianie indeksu. Składanie, identyczność i dodawanie zachowują się stopniowo. Dla homotopii $h$ z~\eqref{cat:eq:cochain-homotopy} definiujemy $\widetilde h^n=(-1)^r h^{n+r}$. Wówczas
\[
 d_{Y[r]}^{n-1}\widetilde h^n+\widetilde h^{n+1}d_{X[r]}^n
 =d_Y^{n+r-1}h^{n+r}+h^{n+r+1}d_X^{n+r}
 =f[r]^n-g[r]^n.
\]
Przesunięcie przechodzi więc na kategorie homotopijne. Wzory na składanie przesunięć wynikają z dodawania indeksów oraz $(-1)^r(-1)^s=(-1)^{r+s}$. Jądro i obraz morfizmu w kategorii abelowej nie zmieniają się jako podobiekty po pomnożeniu go przez $-1$, ponieważ $-\id$ jest izomorfizmem. Stąd obiekty cykli, brzegów i ich ilorazy w $X[r]$ odpowiadają tym z $X$ o indeksie przesuniętym o $r$. Identyfikacje komutują z indukowanymi morfizmami, ponieważ $f[r]^n=f^{n+r}$.
\end{proof}

\subsection{Stożek morfizmu}
\begin{definicja}\label{cat:def:mapping-cone}
Dla morfizmu kompleksów $f:X\to Y$ definiujemy stożek $\operatorname{Cone}(f)$ przez
\[
 \operatorname{Cone}(f)^n=Y^n\oplus X^{n+1},\qquad
 d_{\operatorname{Cone}(f)}^n=
 \begin{pmatrix}d_Y^n&f^{n+1}\\0&-d_X^{n+1}\end{pmatrix}.
\]
Definiujemy $i:Y\to\operatorname{Cone}(f)$ jako włożenie pierwszego składnika i $p:\operatorname{Cone}(f)\to X[1]$ jako projekcję na drugi składnik.
\end{definicja}
\begin{stwierdzenie}\label{cat:prop:cone-complex}
Stożek jest kompleksem, a $i,p$ są morfizmami kompleksów. W kategorii abelowej ciąg
\begin{equation}\label{cat:eq:cone-ses}
 0\to Y\xrightarrow{i}\operatorname{Cone}(f)\xrightarrow{p}X[1]\to0
\end{equation}
jest dokładny i rozkładalny w każdym stopniu. Złożenia kolejnych strzałek $X\xrightarrow f Y\xrightarrow i\operatorname{Cone}(f)\xrightarrow p X[1]\xrightarrow{f[1]}Y[1]$ są zerowe w kategorii homotopijnej.
\end{stwierdzenie}
\begin{proof}
Iloczyn dwóch macierzy różniczek ma zerowe składniki diagonalne przez $d_X^2=d_Y^2=0$. Składnik górny prawy to $d_Y^{n+1}f^{n+1}-f^{n+2}d_X^{n+1}=0$, ponieważ $f$ jest morfizmem kompleksów. Włożenie i projekcja komutują z różniczkami przez postać macierzy, a dokładność stopniowa wynika z biproduktu. Mamy $pi=0$ już w kategorii kompleksów. Dla $if$ wybieramy homotopię $h^n:X^n\to\operatorname{Cone}(f)^{n-1}$, $h^n(x)=(0,x)$, gdzie zapis współrzędnych oznacza włożenie do biproduktu (nie zakładamy istnienia elementów obiektów). Rachunek macierzowy daje $d_{\operatorname{Cone}(f)}h+hd_X=(f,0)^{\mathsf T}=if$. Dla $f[1]p$ wybieramy homotopię $k^n:\operatorname{Cone}(f)^n\to Y[1]^{n-1}=Y^n$ równą projekcji na pierwszy składnik. Wówczas
\[
 d_{Y[1]}^{n-1}k^n+k^{n+1}d_{\operatorname{Cone}(f)}^n
 =[-d_Y^n,0]+[d_Y^n,f^{n+1}]=[0,f^{n+1}].
\]
To jest $f[1]^np^n$, więc wszystkie wskazane złożenia są zerowe w $\cat K(\mathcal A)$.
\end{proof}
\begin{stwierdzenie}\label{cat:prop:cone-homotopic}
Jeżeli $f,g:X\to Y$ są homotopijne, to istnieje izomorfizm kompleksów $\operatorname{Cone}(f)\to\operatorname{Cone}(g)$ komutujący z włożeniami $Y$ i projekcjami na $X[1]$.
\end{stwierdzenie}
\begin{proof}
Dla $f-g=d_Yh+hd_X$ wybieramy w stopniu $n$ macierz
\[
 T^n=\begin{pmatrix}1&h^{n+1}\\0&1\end{pmatrix}.
\]
Górne prawe składniki $d_{\operatorname{Cone}(g)}^nT^n$ i $T^{n+1}d_{\operatorname{Cone}(f)}^n$ to odpowiednio $d_Y^nh^{n+1}+g^{n+1}$ i $f^{n+1}-h^{n+2}d_X^{n+1}$. Są równe przez założenie. Pozostałe składniki są identyczne. Macierz odwrotna powstaje przez zastąpienie $h^{n+1}$ przez $-h^{n+1}$, a ponieważ $T$ komutuje z różniczkami, jego odwrotność również z nimi komutuje. Równania dotyczące włożenia i projekcji wynikają z postaci trójkątnej macierzy.
\end{proof}
\begin{twierdzenie}[Kohomologia stożka]\label{cat:thm:cone-les}
Jeżeli $\mathcal A$ jest abelowa, to istnieje naturalny długi ciąg dokładny
\[
 \cdots\to H^n(X)\xrightarrow{H^n(f)}H^n(Y)
 \xrightarrow{H^n(i)}H^n(\operatorname{Cone}(f))
 \xrightarrow{H^n(p)}H^{n+1}(X)\to\cdots.
\]
W konsekwencji $f$ indukuje izomorfizmy we wszystkich stopniach kohomologii wtedy i tylko wtedy, gdy $\operatorname{Cone}(f)$ jest acykliczny.
\end{twierdzenie}
\begin{proof}
Stosujemy długi ciąg z twierdzenia~\ref{cat:thm:abelian-homology-les} do~\eqref{cat:eq:cone-ses}. Identyfikujemy $H^n(X[1])$ z $H^{n+1}(X)$ na mocy stwierdzenia~\ref{cat:prop:shift-functor}. Morfizm łączący jest $H^{n+1}(f)$: dla cyklu $x:T\to X^{n+1}$ jego podniesienie do stożka ma składowe $(0,x)$, a różniczka podniesienia ma składowe $(f^{n+1}x,0)$. Konstrukcja morfizmu łączącego w~\eqref{cat:eq:homology-connecting}, po odwróceniu indeksów, daje klasę $f^{n+1}x$. Każdy morfizm do kohomologii można lokalnie podnieść do cyklu przez epi; skracając ten epi, uzyskujemy równość morfizmów w kategorii abelowej. Naturalność względem przemiennych kwadratów morfizmów kompleksów wynika z morfizmu ciągów~\eqref{cat:eq:cone-ses} o środkowej macierzy diagonalnej z odpowiednich morfizmów oraz naturalności długiego ciągu.

Jeżeli każde $H^n(f)$ jest izomorfizmem, dokładność daje $H^n(i)=0$ oraz $\ker H^{n+1}(f)=0$, więc $H^n(p)$ jest jednocześnie mono i zerowy. To wymusza $H^n(\operatorname{Cone}(f))=0$. Odwrotnie, przy zerowej kohomologii stożka dokładność po obu stronach $H^n(f)$ daje jego monomorficzność i epimorficzność. W kategorii abelowej oznacza to izomorfizm.
\end{proof}
\begin{definicja}\label{cat:def:quasi-isomorphism}
Morfizm kompleksów w kategorii abelowej nazywamy quasi-izomorfizmem, jeśli indukuje izomorfizmy $H^n(f)$ dla wszystkich $n\in\mathbb Z$. Własność ta zależy tylko od klasy homotopii na mocy twierdzenia~\ref{cat:thm:abelian-homology-functor}.
\end{definicja}
\begin{uwaga}
Równoważność homotopijna jest quasi-izomorfizmem, ponieważ kohomologia jest funktorem na kategorii homotopijnej. Odwrotność nie zachodzi: morfizm $0\to X$ dla kompleksu z przykładu~\ref{cat:ex:acyclic-not-contractible} jest quasi-izomorfizmem, ale nie jest równoważnością homotopijną przez stwierdzenie~\ref{cat:prop:contractible-zero}.
\end{uwaga}

\subsection{Trójkąty wyróżnione i pierwsze aksjomaty}
\begin{definicja}\label{cat:def:triangle}
Niech $\mathcal T$ będzie kategorią addytywną z ustalonym addytywnym autoizomorfizmem $[1]$ (ogólniej autoequiwalencją z ustaloną odwrotnością z dokładnością do izomorfizmu naturalnego). Trójkątem nazywamy diagram
\[
 X\xrightarrow f Y\xrightarrow u U\xrightarrow v X[1].
\]
Morfizmem trójkątów jest trójka $(a,b,c)$ morfizmów na obiektach $X,Y,U$, dla której $bf=f'a$, $cu=u'b$ oraz $a[1]v=v'c$. Jest to izomorfizm trójkątów wtedy i tylko wtedy, gdy $a,b,c$ są izomorfizmami: odwrotne morfizmy spełniają odpowiednie równania po złożeniu równań wyjściowych z odwrotnościami. Rotacją w przód nazywamy trójkąt
\[
 Y\xrightarrow u U\xrightarrow v X[1]\xrightarrow{-f[1]}Y[1],
\]
a rotacją wstecz trójkąt $U[-1]\xrightarrow{-v[-1]}X\xrightarrow f Y\xrightarrow u U$, gdzie stosujemy ustalone utożsamienia odwrotnych przesunięć.
\end{definicja}
\begin{definicja}\label{cat:def:cone-triangle}
W $\cat{K}(\mathcal A)$ trójkątem wyróżnionym nazywamy każdy trójkąt izomorficzny w kategorii trójkątów do trójkąta stożka
\begin{equation}\label{cat:eq:cone-triangle}
 X\xrightarrow{[f]}Y\xrightarrow{[i]}\operatorname{Cone}(f)
 \xrightarrow{[p]}X[1],
\end{equation}
dla pewnego morfizmu kompleksów $f$. Nawiasy kwadratowe oznaczają tu klasy homotopii morfizmów, w odróżnieniu od przesunięcia obiektu. Stwierdzenie~\ref{cat:prop:cone-homotopic} gwarantuje, że zmiana reprezentanta $[f]$ daje izomorficzny trójkąt stożka z identycznościami na $X,Y$.
\end{definicja}
\begin{stwierdzenie}[Warunki TR1--TR3]\label{cat:prop:tr-123}
Trójkąty wyróżnione w $\cat{K}(\mathcal A)$ spełniają następujące warunki.
\begin{enumerate}[label=\textup{(TR\arabic*)}]
 \item Klasa trójkątów wyróżnionych jest zamknięta na izomorfizmy. Trójkąt $X\xrightarrow{\id}X\to0\to X[1]$ jest wyróżniony. Każdy morfizm $X\to Y$ jest pierwszą strzałką pewnego trójkąta wyróżnionego.
 \item Trójkąt jest wyróżniony wtedy i tylko wtedy, gdy jego rotacja w przód jest wyróżniona.
 \item Dla dwóch trójkątów wyróżnionych $X\xrightarrow f Y\xrightarrow u U\xrightarrow v X[1]$ oraz $X'\xrightarrow{f'}Y'\xrightarrow{u'}U'\xrightarrow{v'}X'[1]$ każdy przemienny kwadrat $bf=f'a$ z morfizmami $a:X\to X'$ i $b:Y\to Y'$ można uzupełnić morfizmem $c:U\to U'$ do morfizmu trójkątów.
\end{enumerate}
Te same warunki zachodzą w $\cat{K}^+(\mathcal A)$, $\cat{K}^-(\mathcal A)$ i $\cat{K}^b(\mathcal A)$.
\end{stwierdzenie}
\begin{proof}
\emph{TR1.} Zamkniętość na izomorfizmy wynika z definicji. Dla stożka identyczności kontrakcję daje rodzina $h^n(x,y)=(0,x)$ z $\operatorname{Cone}(\id_X)^n=X^n\oplus X^{n+1}$ do stopnia $n-1$. Istotnie $dh(x,y)=(x,-d_X^nx)$ oraz $hd(x,y)=(0,d_X^nx+y)$, więc $dh+hd=\id$. Na mocy stwierdzenia~\ref{cat:prop:contractible-zero} stożek jest obiektem zerowym w kategorii homotopijnej; zastąpienie go przez $0$ daje żądany trójkąt. Dla dowolnej klasy morfizmu wybieramy reprezentanta i trójkąt~\eqref{cat:eq:cone-triangle}.

\emph{TR2, rotacja w przód.} Wystarczy rozważyć trójkąt stożka, gdyż rotacja przenosi izomorfizmy trójkątów na izomorfizmy. W rachunkach współrzędne oznaczają składowe morfizmu do biproduktu, a nie elementy obiektu kategorii. Oznaczmy $C=\operatorname{Cone}(f)$ i $B=\operatorname{Cone}(i)$ dla $i:Y\to C$. W stopniu $n$ piszemy
\[
 B^n=Y^n\oplus X^{n+1}\oplus Y^{n+1},\qquad
 d_B(y,x,z)=(d_Yy+fx+z,-d_Xx,-d_Yz).
\]
Pomijamy indeksy różniczek i morfizmów, gdy wynikają jednoznacznie z typów składowych. Morfizmy
\[
 r:B\to X[1],\quad r(y,x,z)=x,\qquad
 s:X[1]\to B,\quad s(x)=(0,x,-fx)
\]
komutują z różniczkami: dla $r$ jest to projekcja na $-d_Xx$, a dla $s$ mamy $d_Bs(x)=(0,-d_Xx,d_Yfx)=s(-d_Xx)$, ponieważ $f$ komutuje z różniczkami. Ponadto $rs=\id$. Rodzina $h^n:B^n\to B^{n-1}$ dana przez $h(y,x,z)=(0,0,y)$ spełnia
\[
 d_Bh+hd_B=(y,0,fx+z)=\id_B-sr.
\]
Stąd $r$ jest izomorfizmem w kategorii homotopijnej. Dla włożenia $j:C\to B$ zachodzi $rj=p$. Dla projekcji $q:B\to Y[1]$, $q(y,x,z)=z$, rodzina $k^n:B^n\to Y[1]^{n-1}$ z $k(y,x,z)=y$ daje
\[
 d_{Y[1]}k+kd_B=-d_Yy+(d_Yy+fx+z)=fx+z=q+f[1]r.
\]
A zatem $[q]=-[f[1]][r]$. Trójka $(\id_Y,\id_C,[r])$ jest izomorfizmem trójkąta stożka $i$ na rotację trójkąta stożka $f$.

\emph{TR2, rotacja wstecz.} Żeby uzyskać implikację odwrotną, pokażemy bezpośrednio zamkniętość na rotację wstecz. Dla $C=\operatorname{Cone}(f)$ bierzemy $B'=\operatorname{Cone}(-p[-1])$; jego składowe i różniczka są
\[
 B'^n=X^n\oplus Y^n\oplus X^{n+1},\qquad
 d_{B'}(x,y,z)=(d_Xx-z,d_Yy+fz,-d_Xz).
\]
Morfizmy $r':B'\to Y$, $r'(x,y,z)=fx+y$, oraz $s':Y\to B'$, $s'(y)=(0,y,0)$ są morfizmami kompleksów: $r'd_{B'}=fd_Xx-fz+d_Yy+fz=d_Yr'$ oraz $d_{B'}s'=s'd_Y$. Mamy $r's'=\id_Y$. Dla $h'^n:B'^n\to B'^{n-1}$, $h'(x,y,z)=(0,0,-x)$, rachunek daje
\[
 d_{B'}h'+h'd_{B'}=(x,-fx,z)=\id_{B'}-s'r'.
\]
Zatem $[r']$ jest izomorfizmem. Włożenie $j':X\to B'$ spełnia $r'j'=f$. Projekcja $q':B'\to C$ spełnia $q'(x,y,z)=(y,z)$. Różnica $q'-ir'$ ma składowe $(-fx,z)$ i jest homotopijna z zerem przez $k'^n:B'^n\to C^{n-1}$ dane wzorem $k'(x,y,z)=(0,-x)$:
\[
 d_Ck'+k'd_{B'}=(-fx,d_Xx)+(0,-d_Xx+z)=(-fx,z).
\]
Trójka $(\id_{C[-1]},\id_X,[r'])$ daje więc izomorfizm trójkąta stożka $-p[-1]$ na rotację wstecz trójkąta stożka $f$. Rotacje w przód i wstecz są wzajemnie odwrotne po identyfikacji przesunięć; zastosowanie rotacji wstecz do wyróżnionej rotacji w przód dowodzi drugiej implikacji.

\emph{TR3.} Przenosząc kwadrat przez izomorfizmy trójkątów, sprowadzamy się do trójkątów stożków $f,f'$. Wybieramy reprezentanty $a,b$ w kategorii kompleksów. Przemienność kwadratu w kategorii homotopijnej oznacza istnienie rodziny $h$ z $bf-f'a=d_{Y'}h+hd_X$. Definiujemy morfizm
\[
 c^n:\operatorname{Cone}(f)^n\to\operatorname{Cone}(f')^n,\qquad
 c^n=\begin{pmatrix}b^n&h^{n+1}\\0&a^{n+1}\end{pmatrix}.
\]
Składniki diagonalne równości $d_{\operatorname{Cone}(f')}^nc^n=c^{n+1}d_{\operatorname{Cone}(f)}^n$ są równe, gdyż $a,b$ komutują z różniczkami; dolne lewe są zerowe. Równość górnych prawych składników ma postać
\[
 d_{Y'}^nh^{n+1}+f'{}^{n+1}a^{n+1}
 =b^{n+1}f^{n+1}-h^{n+2}d_X^{n+1},
\]
co jest równaniem homotopii. Mamy $ci=i'b$ i $p'c=a[1]p$ już w kategorii kompleksów, więc $([a],[b],[c])$ jest morfizmem trójkątów. Po przeniesieniu z powrotem przez izomorfizmy dostajemy żądane uzupełnienie.

Wszystkie użyte konstrukcje (przesunięcia, skończone biprodukty i stożki) zachowują ograniczenie z dołu, z góry lub z obu stron; podkategorie są pełne, więc dowód działa bez zmian.
\end{proof}
\begin{uwaga}
Warunki TR1--TR3 nie są jeszcze pełną definicją kategorii triangulowanej. Potrzebny jest dodatkowo aksjomat ośmiościanu, który wiąże trójkąty odpowiadające dwóm składanym morfizmom i ich złożeniu. Uzupełnienie morfizmu trójkątów w TR3 nie jest na ogół jednoznaczne; dowód wymaga wyboru homotopii przemienności kwadratu, a nie tylko klas jego strzałek.
\end{uwaga}
